\section{Future Work}

Several future directions arise naturally from the findings of this thesis and the areas noted in the conclusion of the initial study.

A primary direction is the expansion to larger and more diverse datasets. Because the current dataset was collected via a mobile application, it may underrepresent certain demographic or clinical groups. Incorporating data from clinical records, telemedicine transcripts, or more heterogeneous populations would improve robustness, reduce bias, and strengthen generalizability. Previous research also notes the importance of validating model performance in different populations and evolving epidemic contexts, especially as COVID-19 symptoms and variants change over time.

Another direction involves richer serialization strategies. This study compared list- and text-based formats, but additional formats: structured JSON-like representations, hybrid templates, or templates that emphasize temporal symptom patterns---may further improve consistency and interpretability. Future systems may also integrate multi-turn reasoning to refine predictions over multiple rounds.

Interpretability is another area of development. While attention maps and generated rationales offer intuitive insights, more rigorous interpretability frameworks could strengthen clinician trust. Methods such as attention-based attribution or prompting the model for structured explanations may help reveal how individual features contribute to risk predictions.

Several deployment-oriented opportunities also exist. Evaluating system behavior under real-world usage conditions, including latency, error recovery, and ambiguous user responses---will be essential before clinical or public-health deployment. Mechanisms that detect uncertain predictions, escalate high-risk responses, or provide safety disclaimers would improve the system’s reliability. As noted previously, mobile-based digital tools can play a meaningful role during outbreaks, but ensuring safety and clarity is critical.

Finally, the framework may be extended to other diseases and assessment scenarios. The approach of serializing structured features, fine-tuning LLMs in a few-shot setting, and integrating predictions with conversational explanations is general and not limited to COVID-19. Risk assessment for influenza, chronic conditions, or triage-oriented symptom evaluation are all viable extensions, provided appropriate datasets and validation procedures are available.

In summary, future work should focus on expanding datasets, enhancing serialization and interpretability, validating system behavior in broader contexts, and generalizing the methodology to other areas of digital health. These directions build directly on these findings while supporting stronger, safer, and more versatile health-assessment applications.
